% TEMPLATE-MILC-v3.tex - plantilla maestra UNICA de las guias MILC v3.
%
% La usan las cuatro familias de guias, cada una con su driver:
%   · Grados 6-11          -> scripts/build-guias-g11.py
%   · Bebras               -> scripts/build-guias-bebras.py
%   · Semillero            -> scripts/build-guias-semillero.py
%   · Territorio interior  -> scripts/build-guias-territorio-interior.py
%
% Antes habia tres copias casi identicas (~400 lineas iguales, 6 distintas):
% cada mejora habia que aplicarla tres veces, y en la practica no se hacia
% (el motor de recursos vivia solo en dos de ellas). Lo que varia entre
% familias esta parametrizado en 6 placeholders que cada driver llena
% (se nombran aqui SIN su sintaxis real: el builder reemplaza por texto plano,
% tambien dentro de los comentarios, y un valor multilinea rompe el preambulo):
%
%   HEADER_IZQ        encabezado izquierdo de pagina
%   HEADER_DER        encabezado derecho de pagina
%   PORTADA_ETIQUETA  rotulo sobre el numero grande (GRADO/MOMENTO/MODULO)
%   PORTADA_SERIE     linea de serie de la portada
%   PORTADA_ID        identificador de la guia en la portada
%   PORTADA_CIRCULO   el circulito de grado (°), o vacio si no aplica
%   PDF_TITULO        titulo en texto plano (sin \textbf ni comillas TeX) para
%                     los metadatos del PDF (/Title); lo produce lib_xelatex.texto_pdf
%
% Composicion (auditoria 2026-09-03): las bandas de estacion piden espacio
% (\Needspace*) y prohiben el salto justo despues (\nopagebreak), asi no quedan
% huerfanas al pie; las cajas largas (softbox, drawbox, infoband, puentebox) son
% `breakable`, asi una caja de media pagina ya no arrastra todo a la siguiente
% dejando la anterior al 40 %. Todo texto sobre mostaza o verde va en milcNegro
% (blanco sobre #E5B400 da 1,93:1 y sobre #58A923 2,95:1; ninguno pasa WCAG AA).
% El resto de claves las documenta content/guias/_SCHEMA.md. El script valida
% que no queden placeholders sin reemplazar antes de escribir el archivo final.

% Etiquetado PDF (latex-lab): /Lang, estructura y texto alternativo de las
% figuras, para lectores de pantalla. Debe ir ANTES de \documentclass.
\DocumentMetadata{lang=es-CO,tagging=on}
\documentclass[11pt,letterpaper]{article}
% headheight=24pt: la cabecera derecha lleva dos lineas (identidad de la guia
% y estacion actual). headsep se acorta en la misma medida para que el bloque
% de texto quede exactamente donde estaba con headheight=12pt.
\usepackage[margin=2.0cm,headheight=24pt,headsep=13pt]{geometry}
\usepackage[table]{xcolor}
\usepackage{fontspec}
\usepackage[spanish]{babel}
\usepackage{tikz}
% Librerías TikZ para los diagramas de las guías (bloque `recursos:`):
%  · babel  → babel en español vuelve activos caracteres como «>», y TikZ los usa
%             en sus opciones (>=stealth, ->). Sin esta librería, cualquier
%             diagrama con flechas revienta con «\language@active@arg> has an
%             extra }». Debe ir ANTES de que se dibuje nada.
%  · positioning        → sintaxis «below left=2pt and 3pt».
%  · decorations.pathreplacing → llaves ({brace}) para señalar rangos.
\usetikzlibrary{babel,positioning,decorations.pathreplacing}
\usepackage{graphicx}
% Circuitos: el trabajo de electrónica se hace hoy con micro:bit y sus sensores
% integrados (MakeCode, sin cableado), así que no se necesitan esquemáticos.
% Si algún día entran montajes con protoboard, basta descomentar esta línea:
% los diagramas .tex/.tikz declarados en `recursos:` ya se incrustan solos.
% \usepackage{circuitikz}
\usepackage[most]{tcolorbox}
\usepackage{tabularx}
\usepackage{array}
\usepackage{fancyhdr}
\usepackage{titlesec}
\usepackage[document]{ragged2e}  % bandera a la derecha en todo el cuerpo
\usepackage{enumitem}
\usepackage{microtype}
\usepackage{etoolbox}
\usepackage{needspace}
% hyperref al final del preambulo: metadatos del PDF (titulo, autor, idioma,
% familia), marcadores por estacion (\pdfbookmark) y enlaces sin recuadro.
\usepackage[hidelinks,
  pdftitle={Escuchar sin arreglar: el silencio no es una falla de la conversación},
  pdfauthor={PhD. Álvaro Cárdenas Orozco},
  pdfsubject={Territorio interior · Educación socioemocional · Grado 7° · Momento 1 · MILC v3},
  pdflang={es-CO},
  bookmarksopen=true,bookmarksnumbered=false]{hyperref}

% Helvetica Neue no trae la flecha U+2192, asi que cada «→» escrito literalmente
% en el YAML se compilaba a NADA, en silencio (462 flechas en 61 guias: rutas del
% tipo «paleta → tipografia → lockup» salian rotas). Mapeamos el caracter a su
% version matematica, que si tiene glifo. Asi el autor puede seguir escribiendo
% «→» comodamente en el YAML.
\usepackage{newunicodechar}
\newunicodechar{→}{\ensuremath{\rightarrow}}
% Mismo problema con los simbolos de marca y vinetas: el «✓» de «una palomita ✓»
% salia como cuadrito vacio en el PDF de todas las guias que lo usaban.
\usepackage{pifont}
\newunicodechar{✓}{\ding{51}}
\newunicodechar{✗}{\ding{55}}
\newunicodechar{✔}{\ding{52}}
\newunicodechar{•}{\textbullet}
\newunicodechar{≥}{\ensuremath{\geq}}
\newunicodechar{≤}{\ensuremath{\leq}}

\setmainfont{Helvetica Neue}
\setsansfont{Helvetica Neue}
\setlength{\parindent}{0pt}
\setlength{\parskip}{6pt}
% Sin particion de palabras: en 6.º-7.º una palabra cortada por guion al final
% de linea («Comuni-cacion») frena la lectura mucho mas que una linea corta.
\hyphenpenalty=10000
\exhyphenpenalty=10000
\pretolerance=2000
\tolerance=3000
\setlength{\emergencystretch}{3em}
\linespread{1.10}
\renewcommand{\arraystretch}{1.25}
\setlist{leftmargin=5mm,itemsep=1mm,topsep=1mm}
\newcolumntype{Y}{>{\RaggedRight\arraybackslash}X}

\definecolor{milcMagenta}{HTML}{E6007E}
\definecolor{milcVino}{HTML}{5A0038}
\definecolor{milcTurquesa}{HTML}{0093A5}
\definecolor{milcVerde}{HTML}{58A923}
\definecolor{milcMostaza}{HTML}{E5B400}
\definecolor{milcOcre}{HTML}{B86600}
\definecolor{milcNegro}{HTML}{241816}
\definecolor{milcGris}{HTML}{F7F7F4}
\definecolor{milcLinea}{HTML}{E8E2DD}
\definecolor{milcGradePrimary}{HTML}{0D9488}
\definecolor{milcGradeDark}{HTML}{115E59}
\definecolor{milcGradeSoft}{HTML}{CCFBF1}
\definecolor{milcGradeAccent}{HTML}{CFE7FF}
\definecolor{milcGradeText}{HTML}{FFFFFF}
\color{milcNegro}

\pagestyle{fancy}
\fancyhf{}
% Cabecera de dos lineas: identidad de la guia arriba y, a la derecha y en
% gris, la estacion en curso (\markright desde \milcestacion). Antes de la
% primera estacion la segunda linea va vacia; el \strut mantiene la altura
% para que la linea de arriba no baile entre paginas.
\lhead{\small Territorio interior · Educación socioemocional\\[-1pt]{\footnotesize\strut}}
\rhead{\small Grado 7° · Momento 1 · MILC v3\\[-1pt]{\footnotesize\strut\color{milcNegro!65}\rightmark}}
\cfoot{\small\thepage}
\renewcommand{\headrulewidth}{0pt}

% Sobre mostaza (#E5B400) y verde (#58A923) el texto va en milcNegro; sobre el
% resto de la paleta, en blanco. Se usa dentro de fonttitle/fontupper, donde un
% \color posterior manda sobre coltitle.
\newcommand{\milctintasobre}[1]{%
  \ifboolexpr{test{\ifstrequal{#1}{milcMostaza}} or test{\ifstrequal{#1}{milcVerde}}}%
    {\color{milcNegro}}{\color{white}}}

\newtcolorbox{titlebox}[1]{
  enhanced,colback=#1,colframe=#1,arc=7mm,boxrule=0pt,
  left=7mm,right=7mm,top=3mm,bottom=3mm,
  before skip=8pt,after skip=8pt,
  fontupper=\sffamily\bfseries\Large\color{white}
}

\newtcolorbox{softbox}[3]{
  enhanced,breakable,pad at break=3mm,
  colback=#2,colframe=#1,borderline west={4pt}{0pt}{#1},
  arc=2mm,boxrule=.4pt,left=4mm,right=4mm,top=3mm,bottom=3mm,
  before skip=6pt,after skip=8pt,title={#3},coltitle=white,
  colbacktitle=#1,fonttitle=\sffamily\bfseries\milctintasobre{#1},fontupper=\RaggedRight,
  attach boxed title to top left={xshift=3mm,yshift=-2mm},
  boxed title style={arc=2mm, boxrule=0pt}
}

\newtcolorbox{drawbox}[2]{
  enhanced,breakable,pad at break=4mm,
  colback=white,colframe=#1,arc=4mm,boxrule=1pt,
  left=5mm,right=5mm,top=4mm,bottom=4mm,before skip=8pt,after skip=8pt,
  title={#2},coltitle=white,colbacktitle=#1,fonttitle=\sffamily\bfseries\milctintasobre{#1},
  fontupper=\RaggedRight,
  attach boxed title to top left={xshift=4mm,yshift=-2mm},
  boxed title style={arc=3mm, boxrule=0pt}
}

\newtcolorbox{infoband}[1]{
  enhanced,breakable,pad at break=2mm,colback=#1!8,colframe=#1!50,arc=2mm,boxrule=.5pt,
  left=4mm,right=4mm,top=2mm,bottom=2mm,before skip=4pt,after skip=4pt,
  fontupper=\small\RaggedRight
}

% Puente narrativo entre secciones (contrato editorial Conéctate).
% Da continuidad: "Acabas de X. Ahora vas a Y porque Z."
% Visualmente: banda izquierda en color del grado, texto en itálica pequeña.
\newtcolorbox{puentebox}{
  enhanced,breakable,pad at break=2mm,colback=milcGradeSoft,colframe=milcGradeDark,
  borderline west={3pt}{0pt}{milcGradeDark},
  arc=0mm,boxrule=0pt,left=5mm,right=4mm,top=2.5mm,bottom=2.5mm,
  before skip=4pt,after skip=8pt,
  fontupper=\itshape\small\color{milcNegro}\RaggedRight
}
% La flecha va en modo matematico a proposito: Helvetica Neue NO tiene el glifo
% U+2192, asi que \textrightarrow se compilaba a NADA («Missing character: There
% is no → in font Helvetica Neue») y el puente salia sin flecha. En modo
% matematico la flecha viene de la fuente matematica y si se dibuja.
\newcommand{\puente}[1]{%
  \begin{puentebox}%
    \textcolor{milcGradeDark}{$\rightarrow$}\hspace{2mm}#1%
  \end{puentebox}%
}

\newcommand{\linead}[1]{\noindent\color{milcLinea}\rule{#1}{0.5pt}\color{milcNegro}}
\newcommand{\respuesta}[1]{\par\vspace{2mm}\foreach \n in {1,...,#1}{\noindent\color{milcLinea}\rule{\linewidth}{0.5pt}\par\vspace{5mm}}\color{milcNegro}}
\newcommand{\checkbox}{\textcolor{milcGradeDark}{\fbox{\phantom{x}}}\hspace{5pt}}

% ─── Listas aireadas para las guías socioemocionales ────────────────────
% Componer las verificaciones como «\checkbox texto\\» deja el texto que salta
% de línea debajo del cuadrito y apelmaza el bloque. Estos entornos alinean la
% continuación y airean los ítems. Los usa Territorio Interior; cualquier
% familia puede hacerlo.
\newlist{checklista}{itemize}{1}
\setlist[checklista]{
  label=\textcolor{milcGradeDark}{\fbox{\phantom{x}}},
  leftmargin=9mm, labelsep=3.2mm, itemsep=2.4mm,
  topsep=2.5mm, parsep=0pt, partopsep=0pt, before=\RaggedRight
}
\newlist{pasoslista}{enumerate}{1}
\setlist[pasoslista]{
  label=\textcolor{milcGradeDark}{\sffamily\bfseries\arabic*.},
  leftmargin=8mm, labelsep=3mm, itemsep=2.8mm,
  topsep=1mm, parsep=0pt, partopsep=0pt, before=\RaggedRight
}

\newcommand{\phasecard}[4]{
\begin{tcolorbox}[enhanced,colback=#3,colframe=#2,arc=3mm,boxrule=.5pt,height=3.05cm,valign=center,halign=center,left=1.5mm,right=1.5mm,top=2mm,bottom=2mm]
{\sffamily\bfseries\fontsize{22}{24}\selectfont\textcolor{#2}{#1}}\par
{\sffamily\bfseries\small #4}
\end{tcolorbox}
}

% ── Piezas del rediseno 2026 (legibilidad 6.º-11.º) ───────────────────
% Banda de estacion: reemplaza el titlebox macizo. Titulo a la izquierda,
% dato practico (tiempo/modalidad) a la derecha, en una sola linea.
% Nunca huerfana: pide 9 lineas libres (o salta de pagina rellenando con
% \vfill) y prohibe el salto justo despues de la banda. Nueve y no seis:
% la caja partible que sigue necesita ~80pt (titulo adosado, relleno y dos
% lineas) para arrancar en la misma pagina; con menos, tcolorbox la manda
% entera a la siguiente y la banda se queda sola. El penalty 10000 va
% en `after`, ANTES del espacio posterior: un `after skip` (aunque sea 0pt)
% es un \vskip tras la caja, y ese glue es un punto de corte legal que un
% \nopagebreak escrito despues de \end{tcolorbox} ya no cierra. Cada banda
% deja marcador PDF y actualiza la cabecera derecha.
\newcounter{milcestacion}
\newcommand{\milcestacion}[2]{%
\Needspace*{9\baselineskip}%
\stepcounter{milcestacion}%
\pdfbookmark[1]{#2}{milcestacion\themilcestacion}%
\markright{#2}%
\begin{tcolorbox}[enhanced,colback=#1,colframe=#1,arc=2mm,boxrule=0pt,
  left=5mm,right=5mm,top=2.6mm,bottom=2.6mm,before skip=11pt,
  after={\par\nopagebreak\vskip7pt}]
{\sffamily\bfseries\fontsize{15}{18}\selectfont\milctintasobre{#1}#2}
\end{tcolorbox}}

% Tarjeta de paso numerada: sustituye las tablas «Pilar / Aplicacion», que
% eran formato de consulta docente, no de estudiante. El numero va en un
% circulo sobre el borde para que la secuencia se lea de un vistazo.
\newcommand{\milcpaso}[3]{%
\Needspace{4\baselineskip}%
\begin{tcolorbox}[enhanced,colback=white,colframe=milcLinea,arc=1.5mm,boxrule=.9pt,
  left=14mm,right=4mm,top=3mm,bottom=3mm,before skip=4pt,after skip=4pt,
  fontupper=\RaggedRight,
  overlay={\node[anchor=north west,circle,fill=#2,text=white,inner sep=0pt,
    minimum size=7.4mm,font=\sffamily\bfseries\small]
    at ([xshift=3.6mm,yshift=-3.2mm]frame.north west) {#1};}]
#3
\end{tcolorbox}}

% Casilla marcable de tamano util con lapiz (la anterior era \fbox{\phantom{x}},
% de alto variable segun el texto que la seguia).
\newcommand{\milcmarca}{\framebox[3.8mm]{\rule{0pt}{3.4mm}}}

% Fila marcable de lista suelta (checks de la estacion 1).
\newcommand{\milccheckrow}[1]{%
\par\vspace{1.8mm}\noindent
\makebox[6.4mm][l]{\raisebox{-.55mm}{\textcolor{milcNegro}{\milcmarca}}}%
\parbox[t]{\dimexpr\linewidth-6.4mm\relax}{\RaggedRight #1}\par}

% Celda marcable dentro de la rubrica (tabla de autoevaluacion). Misma
% sangria francesa, pero pensada para vivir dentro de una columna X.
\newcommand{\milcopcion}[1]{%
\makebox[5.2mm][l]{\raisebox{-.55mm}{\textcolor{milcNegro}{\milcmarca}}}%
\parbox[t]{\dimexpr\linewidth-5.2mm\relax}{\RaggedRight #1}}

% ── Recursos visuales de la guía ──────────────────────────────────────
% Declarados en el bloque `recursos:` del YAML y emitidos por el builder.
% \guiaFigura   → imagen rasterizada o PDF (foto, carta celeste, montaje).
% \guiaDiagrama → fuente TikZ/circuitikz (.tex) incrustada como vector:
%                 es la vía para circuitos y esquemáticos editables.
% [#1] = texto alternativo (alt) · #2 = ruta relativa al .tex · #3 = pie
% El alt llega al PDF etiquetado (/Alt) y lo lee el lector de pantalla.
\newcommand{\guiaFigura}[3][]{%
\begin{center}
\includegraphics[width=0.88\linewidth,height=0.38\textheight,keepaspectratio,alt={#1}]{#2}\\[1.5mm]
{\footnotesize\itshape\textcolor{milcNegro!75}{#3}}
\end{center}
}

\newcommand{\guiaDiagrama}[2]{%
\begin{center}
\input{#1}\\[1.5mm]
{\footnotesize\itshape\textcolor{milcNegro!75}{#2}}
\end{center}
}

\begin{document}
\thispagestyle{empty}

% ── Portada ───────────────────────────────────────────────────────────
% Antes: un rectangulo de color a pagina completa. Se veia bien en pantalla,
% pero en la fotocopiadora del colegio se comia el toner y salia gris sucio.
% Ahora la banda ocupa el tercio superior y el resto va en blanco.
\begin{tikzpicture}[remember picture,overlay]
  \path[fill=milcGradePrimary]
    (current page.north west) rectangle ([yshift=-10.6cm]current page.north east);

  \node[anchor=north west,text=milcGradeText] at ([xshift=2.0cm,yshift=-1.55cm]current page.north west)
    {\sffamily\bfseries\fontsize{12}{14}\selectfont MOMENTO};
  \node[anchor=north west,text=milcGradeText,opacity=.85] at ([xshift=2.0cm,yshift=-2.35cm]current page.north west)
    {\sffamily\bfseries\fontsize{8.5}{10}\selectfont TERRITORIO INTERIOR · SOCIOEMOCIONAL};
  \node[anchor=north east,text=milcGradeText,opacity=.85,align=right] at ([xshift=-2.0cm,yshift=-1.55cm]current page.north east)
    {\sffamily\bfseries\fontsize{9}{12}\selectfont I.E. Sor María Juliana\\Cartago, Valle del Cauca};

  \node[anchor=west,text=milcGradeText] (gradonum) at ([xshift=1.85cm,yshift=-7.1cm]current page.north west)
    {\sffamily\bfseries\fontsize{112}{112}\selectfont 1};
  % El circulito de grado (°) solo aplica a las guias de grado («11°»); en
  % Bebras y Semillero el numero grande es un momento/modulo, no un grado.
  % Se posiciona relativo al nodo `gradonum`, no en coordenadas absolutas.
  

  \node[anchor=west,text=milcGradeText,text width=11.4cm,align=left] at ([xshift=1.5cm,yshift=0.5cm]gradonum.east)
    {
     {\sffamily\bfseries\fontsize{9}{11}\selectfont\MakeUppercase{Territorio interior · Socioemocional}}\\[3mm]
     {\sffamily\bfseries\fontsize{21}{25}\selectfont\setlength{\baselineskip}{27pt}Escuchar\\[4pt]sin arreglar\\[4pt]el silencio no es una falla}\\[3.5mm]
     {\sffamily\bfseries\fontsize{15}{18}\selectfont Grado 7° · Momento 1}
    };
\end{tikzpicture}

\vspace*{9.4cm}
\vfill

{\sffamily\bfseries\fontsize{8}{10}\selectfont\textcolor{milcGradeDark}{LO QUE TE LLEVAS HOY}}

\begin{tcolorbox}[enhanced,colback=white,colframe=milcNegro,arc=2mm,boxrule=1.6pt,
  left=5mm,right=5mm,top=4mm,bottom=4mm,before skip=3pt,after skip=12pt,
  fontupper=\RaggedRight\fontsize{11}{15.5}\selectfont]
Tu registro de escucha: una conversación real de tres minutos sostenida sin dar un solo consejo —con las tres preguntas que usaste en vez de soluciones, el resumen que le devolviste a la otra persona y lo que te costó callar.
\end{tcolorbox}

\begin{tcolorbox}[enhanced,colback=milcGris,colframe=milcGris,arc=2mm,boxrule=0pt,
  left=5mm,right=5mm,top=3.5mm,bottom=3.5mm,after skip=12pt]
{\sffamily\bfseries\fontsize{8}{10}\selectfont\textcolor{milcNegro!55}{ESTA GUÍA ES DE}}\\[3mm]
\begin{tabularx}{\linewidth}{X p{3.2cm} p{3.2cm}}
\linead{\linewidth} & \linead{3.0cm} & \linead{3.0cm}\\[-1mm]
{\fontsize{8.5}{10}\selectfont\textcolor{milcNegro!55}{Nombre completo}} &
{\fontsize{8.5}{10}\selectfont\textcolor{milcNegro!55}{Grupo}} &
{\fontsize{8.5}{10}\selectfont\textcolor{milcNegro!55}{Fecha}}\\
\end{tabularx}
\end{tcolorbox}

\vfill

\noindent\textcolor{milcLinea}{\rule{\linewidth}{1pt}}\\[2mm]
{\sffamily\bfseries\fontsize{9.5}{12}\selectfont Autor: Dr. Álvaro Cárdenas Orozco}\\[1mm]
{\sffamily\fontsize{8.5}{11}\selectfont\textcolor{milcNegro!55}{Metodología MILC · Apertura ancestral · Escucha activa · Triángulo Dussel-Estoicismo-Floridi}}

\newpage

\section*{Apertura: Escuchar sin arreglar: el silencio no es una falla de la conversación}

\begin{softbox}{milcGradeDark}{milcGradeSoft}{Saber ancestral}
En los pueblos indígenas de Colombia hay una forma de reunirse que no se parece ni a una clase ni a una asamblea de colegio: el \textbf{círculo de la palabra}. Se sientan \textbf{en círculo} ---sin tarima, sin cabecera, sin nadie al frente---, lo conducen \textbf{los mayores}, y la palabra \textbf{circula}: cada quien habla en su turno y \textbf{termina de hablar}. La ONIC y el CRIC los siguen convocando hoy, y hay círculos documentados que duran \textbf{cinco horas}. \textbf{Tres rasgos nos importan.} \textbf{Primero: la palabra circula, no compite.} No hay que ganarse el turno levantando más la voz ni interrumpiendo; el turno llega. \textbf{Segundo: el que conduce es el que más escucha.} Los mayores no dirigen hablando: dirigen sosteniendo el espacio para que hablen los demás. \textbf{Tercero: el círculo no se acaba cuando alguien encuentra la respuesta,} sino cuando todos alcanzaron a decir lo suyo ---y por eso hay que ir con tiempo---. \emph{Compáralo con una conversación de recreo, donde el que habla apenas alcanza a decir tres frases antes de que alguien lo tape con un «eso no es nada, a mí me pasó\ldots». No es que allá tengan más paciencia: es que saben que la palabra ajena necesita espacio para salir completa.}
\end{softbox}

\begin{softbox}{milcTurquesa}{milcGris}{Saber contemporáneo}
¿Y sirve de algo escuchar, si uno no aporta nada útil? \textbf{Se midió.} \textbf{Guy Itzchakov}, \textbf{Avraham Kluger} y \textbf{Dotan Castro} hicieron cuatro experimentos: unas personas hablaban de un tema que les importaba frente a alguien que escuchaba \textbf{bien} ---con atención, sin juzgar y sin apurar---, y otras frente a alguien distraído. \textbf{Los bien escuchados bajaron su ansiedad}, y eso produjo algo que nadie esperaría: \textbf{dejaron de defenderse}. Empezaron a ver \textbf{las dos caras} de su propio asunto, reconocieron sus contradicciones ---y, esto es lo notable, \textbf{pudieron tolerarlas sin angustiarse}--- y sus posturas se volvieron menos extremas (Itzchakov, Kluger y Castro, 2017). \textbf{Traducido a tu vida:} el que escucha bien \textbf{no necesita dar un solo consejo} para que el otro piense mejor. \emph{La escucha sola ya lo hace.} Y al revés: cuando alguien se siente juzgado o apurado, se pone a defender su versión ---y quien está defendiéndose no está pensando---.
\end{softbox}

\begin{softbox}{milcMostaza}{milcGris}{Pregunta-puente}
\textit{En el círculo, \textbf{el que conduce es el que más calla}, y nadie interrumpe para acortar camino. \textbf{¿Cuántas veces esta semana tapaste a alguien con un «lo que tienes que hacer es\ldots»} antes de que terminara de contar lo que le pasaba?}
\end{softbox}

\begin{softbox}{milcVerde}{milcGris}{Saber y saber hacer}
Al terminar podrás: (1) \textbf{distinguir} escuchar de \emph{esperar turno}, y reconocer en conversaciones reales cuál de las dos hiciste; (2) \textbf{explicar} por qué dar la solución demasiado pronto cierra la conversación, y por qué el silencio hace trabajo; (3) \textbf{crear} tres preguntas que abren en lugar de resolver; (4) \textbf{sostener} una conversación real de tres minutos sin dar un solo consejo, y devolver un resumen que la otra persona reconozca como suyo.
\end{softbox}



\small
\begin{tabularx}{\textwidth}{>{\bfseries}p{3.0cm}Y}
\rowcolor{milcGradePrimary!12}Autor & Dr. Álvaro Cárdenas Orozco\\
DBA & Comprende los fenómenos que afectan a su territorio, regula sus emociones ante ellos y acompaña a otros con empatía (transversal 6°--11°, articulado con la política «Escuela, territorio de vida» del MEN, Resolución 006519 de 2025, y la Ley 1620 de 2013)\\
Referentes & Primera ayuda psicológica (OMS, 2011/2012) · Directrices IASC (2007) · USGS y Servicio Geológico Colombiano · Pueblo nasa y la minga (CRIC) · Dussel · Estoicismo · Floridi\\
Producto final & Tu registro de escucha: una conversación real de tres minutos sostenida sin dar un solo consejo —con las tres preguntas que usaste en vez de soluciones, el resumen que le devolviste a la otra persona y lo que te costó callar.\\
\end{tabularx}
\normalsize

\puente{Este es el primer momento del año, y va primero por una razón: \textbf{todo lo que viene después} ---herir con la palabra, reparar, volver a confiar--- \textbf{supone que alguien sea capaz de oír al otro hasta el final}.}

\milcestacion{milcGradeDark}{Ruta MILC: así vamos a aprender}
\begin{tabularx}{\textwidth}{>{\centering\arraybackslash}X>{\centering\arraybackslash}X>{\centering\arraybackslash}X>{\centering\arraybackslash}X}
\phasecard{1}{milcVerde}{milcGris}{Escucha\\[-1mm]\small Observo, escucho y pregunto.} &
\phasecard{2}{milcTurquesa}{milcGris}{Sistematización\\[-1mm]\small Callo y entiendo.} &
\phasecard{3}{milcMagenta}{milcGris}{Praxis\\[-1mm]\small Construyo y aplico.} &
\phasecard{4}{milcVino}{milcGris}{Evaluación\\[-1mm]\small Reflexión liberadora.}
\end{tabularx}


\puente{Primero, dos conversaciones tuyas: una en la que te sentiste escuchado y otra en la que no. Vas a encontrar la diferencia solo, antes de que nadie te la explique.}

\milcestacion{milcVerde}{Estación 1 de 3 · Escucha}

\begin{softbox}{milcVerde}{milcGris}{Escena para escuchar}
\textbf{Actividad 1 · IDENTIFICA --- Dos conversaciones.} \textbf{Nadie va a leer esta página.} \textbf{Primera.} Recuerda una vez en que le contaste algo importante a alguien y \textbf{te sentiste escuchado}: ¿quién era, qué hizo con la cara y con el cuerpo, qué te dijo, y ---sobre todo--- \textbf{qué \emph{no} hizo}? \textbf{Segunda.} Ahora una en que contaste algo y \textbf{no} te sentiste escuchado. ¿En qué segundo exacto se te quitaron las ganas de seguir contando? \emph{Fíjate si fue cuando te dieron un consejo, cuando te contaron un caso peor o cuando alguien miró el celular.} \textbf{Tercera parte, la incómoda.} Recuerda una vez en que \textbf{tú} fuiste el que no escuchó: alguien te estaba contando algo y tú lo arreglaste, lo apuraste o lo tapaste con tu propia historia. \emph{Todo el mundo tiene una; el que no la encuentra es porque lo hace muy seguido.}
\end{softbox}

{\sffamily\bfseries\fontsize{8}{10}\selectfont\textcolor{milcNegro!55}{MARCA CUANDO LO HAYAS HECHO}}
\milccheckrow{Escribiste las dos conversaciones y, en la primera, \textbf{qué no hizo} la persona que te escuchó bien.}
\milccheckrow{Identificaste el segundo exacto en que se te quitaron las ganas de seguir contando.}
\milccheckrow{Anotaste una vez en que el que no escuchó fuiste tú, sin justificarte.}

\begin{infoband}{milcVerde}


\textbf{En tu cuaderno (Actividad 1 · IDENTIFICA).} Título: «Actividad 1. Dos conversaciones». \textbf{Formato:} las dos conversaciones, con el momento en que la segunda se cerró + tu propio caso. \textbf{Extensión:} media página. \textbf{Sabes que terminaste cuando} puedes nombrar \textbf{una cosa concreta que hizo} la persona que sí te escuchó. Esta página es tuya: \textbf{nadie más tiene que leerla si tú no quieres}.
\end{infoband}

\puente{Ya las tienes. Ahora, qué es escuchar de verdad, qué hace el silencio y por qué \emph{arreglar} suele ser una manera educada de interrumpir.}

\milcestacion{milcTurquesa}{Fase 2 · Sistematización: entender la escucha}

\begin{softbox}{milcTurquesa}{milcGris}{Cuatro claves sobre escuchar de verdad}
Cuatro claves para que escuchar deje de ser «quedarse callado» y pase a ser algo que se hace a propósito.
\end{softbox}

\milcpaso{1}{milcTurquesa}{\textbf{Escuchar no es esperar turno.} Hay una manera muy común de \emph{parecer} que se escucha: quedarse callado mientras se prepara la respuesta. \emph{Se nota, y se nota rápido:} el que espera turno contesta antes de que el otro termine la última palabra, porque su frase ya estaba lista. \textbf{En el círculo de la palabra eso no cabe:} como el turno llega solo, no hay que pelearlo, y entonces la cabeza queda libre para oír. \textbf{Señal para detectarte:} si mientras el otro habla ya sabes exactamente qué vas a decir, \emph{no estás escuchando: estás redactando}.}
\milcpaso{2}{milcTurquesa}{\textbf{El silencio trabaja, y casi nadie lo deja trabajar.} Después de que alguien dice algo difícil viene una pausa, y esa pausa \textbf{incomoda}: se siente como una falla de la conversación, como si a uno le tocara llenarla. \emph{No es una falla.} Es el momento en que el que habla decide si sigue ---y lo más importante suele venir \textbf{después} de esa pausa, no antes---. \textbf{Regla práctica:} cuando el otro se calle, \textbf{cuenta hasta tres antes de decir nada}. Vas a descubrir que en la mayoría de los casos no te toca hablar a ti: el otro retoma solo.}
\milcpaso{3}{milcTurquesa}{\textbf{Arreglar es una forma educada de interrumpir.} «Lo que tienes que hacer es\ldots», «no te preocupes, eso pasa», «yo que tú\ldots». \emph{Suenan a ayuda y funcionan como puerta cerrada:} le dicen al otro que ya entendiste ---aunque apenas llevas medio minuto--- y que la conversación puede terminar. \textbf{Peor todavía es el caso propio:} «a mí me pasó algo parecido» le quita el tema al que lo trajo. \emph{Y hay algo que el estudio de Itzchakov y su equipo deja claro:} cuando alguien se siente juzgado o apurado, se pone a \textbf{defender su versión}, y quien se defiende \textbf{deja de pensar}. \textbf{Aconsejar temprano no acelera la solución: la aleja.}}
\milcpaso{4}{milcTurquesa}{\textbf{Preguntar en vez de resolver.} Si no vas a dar consejos, ¿qué haces con la boca? \textbf{Preguntas}, pero de un tipo preciso: las que \textbf{abren} en vez de las que cierran. \emph{Cierran} las de sí o no («¿y le dijiste?»), las que ya traen la respuesta adentro («¿no será que estás exagerando?») y las de curiosidad de chisme («¿y quién más estaba?»). \emph{Abren} tres: \textbf{«¿cómo fue?»}, \textbf{«¿qué fue lo más difícil de eso?»} y \textbf{«¿qué necesitas ahora?»}. \textbf{Esa última es la que cambia todo,} porque muchas veces la respuesta no es «una solución» sino «que alguien me oiga», \emph{y hasta que no se pregunta, el que escucha no tiene cómo saberlo}.}

\begin{drawbox}{milcTurquesa}{La conversación que no se arregla, paso por paso}
\begin{pasoslista}
\item \textbf{Guardo el celular} donde no se vea. \emph{Boca abajo sobre la mesa todavía se ve.}
\item \textbf{Miro y no hablo} mientras el otro cuenta. Ni un «ajá» cada tres segundos.
\item \textbf{Cuento hasta tres} cuando se calla. \emph{Casi siempre retoma solo, y ahí viene lo importante.}
\item \textbf{Pregunto para abrir:} «¿cómo fue?», «¿qué fue lo más difícil?».
\item \textbf{Resumo y devuelvo:} «entonces lo que te tiene mal es \_\_\_\_, ¿voy bien?». \emph{Que él corrija.}
\item \textbf{Pregunto qué necesita} ---y solo si pide consejo, doy consejo---.
\end{pasoslista}
\end{drawbox}

\begin{softbox}{milcMostaza}{milcGris}{Errores comunes y su corrección}
\textbf{Esperar turno:} tu respuesta ya estaba lista antes de que terminara. \textbf{El caso propio:} «a mí me pasó algo parecido» le quita el tema al dueño. \textbf{Llenar el silencio:} la pausa era del otro, no tuya. \textbf{El consejo temprano:} cierra la puerta y suena a «ya entendí, cambiemos de tema». \textbf{Minimizar:} «eso no es nada» y «hay gente peor» no consuelan a nadie. \textbf{El interrogatorio:} preguntar mucho tampoco es escuchar si son preguntas de chisme. \textbf{El celular:} mirarlo una sola vez ya le dijo al otro dónde estaba tu atención.
\end{softbox}

\begin{infoband}{milcTurquesa}


\textbf{En tu cuaderno (Actividad 2 · EXPLICA).} Título: «Actividad 2. Abrir y cerrar». \textbf{Formato:} dos columnas ---«preguntas que abren» y «frases que cierran»---, con al menos cuatro en cada una, sacadas de tus propias conversaciones de la Actividad 1. \textbf{Extensión:} media página. \textbf{Sabes que terminaste cuando} puedes explicar por qué \textbf{«no te preocupes» cierra} una conversación en vez de calmarla.
\end{infoband}

\puente{Ya sabes qué es. Es hora de hacerlo: tres minutos callado, con preguntas en vez de soluciones. \emph{Es más difícil de lo que suena.}}

\milcestacion{milcMagenta}{Fase 3 · Praxis: sostener la escucha}

\begin{softbox}{milcMagenta}{milcGris}{Cómo se sostiene una conversación sin arreglarla}
Ahora se practica. Primero en parejas y con cronómetro, porque callar tres minutos es mucho más difícil de lo que parece; después, con alguien de verdad.
\end{softbox}

\milcpaso{1}{milcMagenta}{\textbf{Los tres minutos (6 min, en parejas).} Uno habla de algo que le importa; el otro \textbf{no dice absolutamente nada} durante tres minutos ---puede mirar, asentir, nada más--- y luego cambian. \emph{Al terminar, el que habló responde una sola pregunta: «¿en qué momento sentiste que te estaban oyendo?».} \textbf{Casi siempre pasa lo mismo:} al que calló se le hicieron eternos los tres minutos, y al que habló se le pasaron volando.}
\milcpaso{2}{milcMagenta}{\textbf{Las preguntas que abren (8 min).} Escribe \textbf{tres preguntas} tuyas que sirvan para que alguien siga contando ---no para averiguar detalles---. Pásalas por el filtro: \emph{¿se puede contestar con sí o no?} Entonces no sirve. \emph{¿Ya trae adentro lo que yo creo?} Tampoco. \textbf{Guarda las tres:} son la herramienta que vas a usar el resto del año, y también en el momento 5, cuando toque abrir una conversación difícil.}
\milcpaso{3}{milcMagenta}{\textbf{El resumen que devuelve (8 min, en parejas).} Segunda ronda: ahora el que escucha, al final, \textbf{devuelve un resumen} de dos frases ---«entonces lo que te tiene mal es \_\_\_\_»--- y el que habló dice \textbf{si le quedó bien o no}. \emph{Este paso es el que convierte la escucha en algo verificable:} si el resumen falla, no hubo escucha, por callado que uno haya estado.}
\milcpaso{4}{milcMagenta}{\textbf{La conversación real (fuera de clase).} Escoge a \textbf{una persona} ---de la casa, del salón, del barrio--- y sostén con ella una conversación de \textbf{tres minutos sin dar un solo consejo}. \textbf{No le avises que es una tarea:} el punto es que sea una conversación de verdad. \emph{Y anota después lo que más cuesta anotar: cuántas veces estuviste a punto de interrumpir, y qué ibas a decir.}}

\begin{drawbox}{milcMagenta}{Antes de cerrar tu registro de escucha}
\begin{checklista}
\item Sostuviste tres minutos \textbf{sin dar un solo consejo} y sin contar un caso tuyo.
\item Usaste al menos \textbf{dos de tus tres preguntas} que abren.
\item Le devolviste un \textbf{resumen} y la otra persona lo reconoció como suyo (o lo corrigió).
\item Dejaste al menos \textbf{un silencio} sin llenarlo, y anotaste qué pasó después.
\item Anotaste \textbf{cuántas veces estuviste a punto de interrumpir} y qué ibas a decir.
\item Entendiste que escuchar bien \textbf{no es quedarse callado}: es hacer que el otro alcance a terminar.
\end{checklista}
\end{drawbox}

\begin{softbox}{milcMagenta}{milcGris}{Plantilla de trabajo}
\textbf{Mis tres preguntas que abren.} \emph{1) \_\_\_\_ · 2) \_\_\_\_ · 3) \_\_\_\_} (ninguna se contesta con sí o no). \\[1mm]
\textbf{El resumen que devuelvo.} \emph{«Entonces lo que te tiene \_\_\_\_ es \_\_\_\_, ¿voy bien?»} \\[1mm]
\textbf{La pregunta final, siempre la misma.} \emph{«¿Qué necesitas ahora?»} ---y si la respuesta es «nada, solo contarlo», \textbf{ya cumpliste}---.
\end{softbox}

\begin{infoband}{milcMagenta}


\textbf{En tu cuaderno (Actividad 3 · CREA).} Título: «Actividad 3. Mi registro de escucha». \textbf{Formato:} las tres preguntas + qué contó la persona + el resumen que devolviste y si acertaste + cuántas veces estuviste a punto de interrumpir. \textbf{Extensión:} una página. \textbf{Sabes que terminaste cuando} puedes escribir \textbf{algo que la otra persona dijo después de un silencio} que no llenaste.
\end{infoband}

\puente{El producto es un registro de escucha. \emph{No se evalúa por lo que dijiste, sino por lo que la otra persona alcanzó a decir gracias a ti.}}

\milcestacion{milcMostaza}{Producto de la sesión}
\textbf{Registro de escucha: tres minutos sin consejos, con preguntas que abren y un resumen devuelto}.
\begin{softbox}{milcMostaza}{milcGris}{Criterios de calidad}
(1) La conversación se sostuvo sin dar consejos y sin desviarse hacia un caso propio. (2) Las tres preguntas abren: ninguna se contesta con sí o no, ni trae la respuesta adentro. (3) El resumen devuelto fue reconocido como suyo por la otra persona, o corregido por ella. (4) Se registró al menos un silencio sostenido y lo que la otra persona dijo después. (5) El registro anota los impulsos de interrumpir, no solo los aciertos.
\end{softbox}

\puente{Antes de cerrar, revisa una cosa: si en tus tres minutos apareció la frase «a mí me pasó algo parecido», la conversación dejó de ser del otro.}

\milcestacion{milcVino}{Evaluación liberadora · cinco dimensiones}

\begin{softbox}{milcVino}{milcGris}{Sentido de la evaluación}
Evaluar no es solo revisar una nota. Es reconocer qué aprendí, qué cambió en mí, qué emociones manejé, cómo me relaciono con la ciudad y la comunidad, y qué le dejaré a quien viene detrás.
\end{softbox}

\textbf{Mi reflexión liberadora en cinco dimensiones.} Completa cada frase iniciada:

\small
\begin{tabularx}{\textwidth}{>{\bfseries\RaggedRight\arraybackslash}p{3.4cm}Y>{\RaggedRight\arraybackslash}p{4.6cm}}
\rowcolor{milcVino}\color{white}Dimensión & \color{white}\bfseries Mi respuesta (frame starter) & \color{white}\bfseries Algo que descubrí\\
1. Personal & Lo que más cambió en mí fue \linead{6.5cm} & \linead{4.2cm}\\
2. Emocional & La emoción más fuerte fue \linead{2.5cm} y la regulé \linead{3cm} & \linead{4.2cm}\\
3. Ciudadana & En democracia esto importa porque \linead{6cm} & \linead{4.2cm}\\
4. Local (Cartago) & En mi barrio sería útil para \linead{6.5cm} & \linead{4.2cm}\\
5. Intergeneracional & A mi abuela/abuelo le diría \linead{6.5cm} & \linead{4.2cm}\\
\end{tabularx}
\normalsize

\begin{infoband}{milcVino}
\textbf{Para la próxima sustentación.} Vuelve a esta tabla en seis meses. Verás que las respuestas cambian — no porque lo aprendido cambie, sino porque tú cambias. La evaluación liberadora es una conversación contigo a través del tiempo.
\end{infoband}


\puente{Cerramos con tres voces: la escucha que funda comunidad, la disposición a cambiar de opinión con lo que se oye, y qué se cultiva cuando alguien escucha bien.}

\milcestacion{milcGradeDark}{Triángulo de pensamiento · cierre filosófico}

\begin{softbox}{milcVino}{milcGris}{Dussel · lente del nosotros}
\textit{«Aquel que es capaz de escuchar silenciosamente al Otro como otro es el que puede constituir una comunidad y no una mera sociedad totalizada.»}

{\footnotesize\itshape\color{milcNegro!70}--- Enrique Dussel · Filosofía de la liberación (1977)}

Dussel pone la escucha en un lugar que sorprende: \textbf{no como una buena costumbre, sino como lo que hace la diferencia entre un grupo y un montón de gente junta}. \emph{Un salón donde nadie termina de hablar es una «sociedad totalizada»:} funciona, se organiza, saca notas ---y nadie sabe realmente qué le pasa al de al lado---. \textbf{Y fíjate en el adverbio:} \emph{silenciosamente}. No dice escuchar para responder mejor, ni para ayudar: dice \textbf{escuchar en silencio al otro \emph{como otro}}, es decir, sin convertirlo en una versión de uno mismo.

\textbf{¿En mi salón hay alguien que nunca alcanza a terminar de hablar porque siempre lo tapan?}
\end{softbox}
\linead{\linewidth}\\[1mm]
\linead{\linewidth}

\begin{softbox}{milcMostaza}{milcGris}{Estoicismo · Marco Aurelio · Meditaciones VI, 21 (c. 175 d.C.) · cuidado interior}
\textit{«Si alguien puede refutarme y probar de modo concluyente que pienso o actúo incorrectamente, de buen grado cambiaré de proceder. Pues persigo la verdad, que no dañó nunca a nadie; en cambio, sí se daña el que persiste en su propio engaño e ignorancia.»}

{\footnotesize\itshape\color{milcNegro!70}--- Marco Aurelio · Meditaciones VI, 21 (c. 175 d.C.)}

Esto lo escribe el hombre más poderoso de su mundo, y dice algo incómodo: \textbf{estoy dispuesto a que lo que oiga me cambie}. \emph{Ahí está la diferencia entre oír y escuchar:} el que espera turno ya decidió que no va a mover nada de lo suyo. \textbf{Y el remate es exacto:} el que insiste en su propio error \textbf{se daña a sí mismo}. Escuchar no es un favor que uno le hace al otro ---es la única manera de enterarse de lo que uno tiene mal---. \emph{Por eso el estudio de Itzchakov encaja aquí:} al que \emph{se sintió} escuchado se le aflojaron las certezas, y pudo mirar las dos caras sin angustiarse.

\textbf{¿Cuándo fue la última vez que cambié de opinión por algo que alguien me dijo?}
\end{softbox}
\linead{\linewidth}\\[1mm]
\linead{\linewidth}

\begin{softbox}{milcTurquesa}{milcGris}{Floridi · lente de la infoesfera}
\textit{«El florecimiento de las entidades informacionales y de toda la infoesfera debe promoverse preservando, cultivando y enriqueciendo sus propiedades.»}

{\footnotesize\itshape\color{milcNegro!70}--- Luciano Floridi · La ética de la información (2013; trad.)}

Floridi habla de \textbf{cultivar}, y esa palabra sirve exactamente aquí. \emph{Escuchar bien no es no hacer nada:} es \textbf{hacer que exista algo que sin ti no existiría} ---lo que el otro alcanzó a pensar mientras hablaba---. \textbf{Y hay una versión de esto que se te va a aparecer todo el año:} en un chat nadie deja silencios, todo el mundo responde de una y encima queda escrito. \emph{Lo que se dice apurado se conserva igual que lo que se dice pensado,} y por eso las conversaciones difíciles casi nunca deberían ocurrir por chat.

\textbf{¿Cuántas conversaciones importantes tuve esta semana por chat, que habrían salido mejor de frente?}
\end{softbox}
\linead{\linewidth}\\[1mm]
\linead{\linewidth}

\begin{infoband}{milcNegro}
\textbf{Nota para el docente.} Las tres son citas textuales y llevan su referencia debajo. Puedes remitir a la obra en clase.
\end{infoband}

\milcestacion{milcVino}{Mi autoevaluación · marca una casilla por fila}

{\small
\setlength{\extrarowheight}{2.2pt}
\begin{tabularx}{\textwidth}{>{\RaggedRight\arraybackslash\bfseries}p{3.0cm}YYY}
\rowcolor{milcVino}\color{white}\bfseries Criterio & \color{white}\bfseries Lo logré & \color{white}\bfseries En proceso & \color{white}\bfseries Necesito apoyo\\[.6mm]
Comprensión del tema & \milcopcion{Explico las ideas con mis palabras.} & \milcopcion{Reconozco las ideas, falta precisión.} & \milcopcion{Necesito repasar lo básico.}\\[.8mm]
\rowcolor{milcGris}Calidad de la escucha & \milcopcion{Sostuve la conversación sin dar consejos y la otra persona siguió hablando por su cuenta.} & \milcopcion{Aguanté un rato, pero terminé contando mi propio caso parecido.} & \milcopcion{Interrumpí con la solución en el primer minuto.}\\[.8mm]
Aplicación y producto & \milcopcion{Mi producto cumple los criterios.} & \milcopcion{Está armado, con ajustes pendientes.} & \milcopcion{Aún está en construcción.}\\[.8mm]
\rowcolor{milcGris}Reflexión 5 dimensiones & \milcopcion{Respondí cada una con honestidad.} & \milcopcion{Respondí algunas con detalle.} & \milcopcion{Mis respuestas son muy generales.}\\[.8mm]
Triángulo de pensamiento & \milcopcion{Conecté las 3 voces con mi trabajo.} & \milcopcion{Conecté al menos una.} & \milcopcion{No las relaciono con mi proceso.}\\[.8mm]
\end{tabularx}}

\vspace{3mm}
\textbf{Hoy entendí que:}\\[1mm]
\linead{\linewidth}\\[1mm]
\linead{\linewidth}

\textbf{Mi compromiso para la próxima sesión es:}\\[1mm]
\linead{\linewidth}\\[1mm]
\linead{\linewidth}

\begin{softbox}{milcMostaza}{milcGris}{Cita de despedida · Marco Aurelio}
\textit{``El obstáculo en el camino se vuelve el camino. Lo que estorba al camino se vuelve camino.''}\\[1mm]
Cada error de hoy, bien observado, se convierte en pieza del camino que pisarás mañana.
\end{softbox}

\small
\begin{tabularx}{\textwidth}{>{\bfseries}p{3.6cm}Y}
\rowcolor{milcGradeDark!12}Nombre completo: & \linead{10cm}\\
Fecha: & \linead{4cm} \quad Curso/grupo: \linead{4cm}\\
Mi firma: & \linead{9cm}\\
\end{tabularx}
\normalsize

\begin{infoband}{milcGradeDark}
\textbf{Para la próxima sesión.} Dos cosas. \textbf{Primero, escucha a alguien con quien casi no hablas} ---no a tu mejor amigo, que es lo fácil: al que se sienta solo, al que acaba de llegar, a alguien mayor de tu casa---. Tres minutos, tus tres preguntas, ningún consejo. \emph{Anota una cosa que aprendiste de esa persona y que no tenías cómo saber.} \textbf{Segundo, pregunta en tu casa cómo se hablaba antes:} averigua con alguien mayor \textbf{dónde se conversaba en su época} ---en el corredor, en la cocina, en el andén al atardecer, después del almuerzo del domingo---, \textbf{quién hablaba y quién escuchaba}, y si había algún lugar o momento en que los muchachos podían decir lo suyo. \textbf{Trae:} lo que aprendiste escuchando y la costumbre que te contaron. \emph{Vas a encontrar algo que sorprende: casi todas las familias tenían un lugar y una hora para conversar, y casi ninguna los tiene ahora.}
\end{infoband}

\end{document}
